\section{引言}
中心 Hardy--Littlewood 极大算子的弱 $(1,1)$ 不等式将局部平均的几何结构与水平集的体积联系起来。本文固定 Lebesgue 测度及轴平行方体
\[
 Q(x,r)=x+(-r,r)^n,\qquad
 M_nf(x)=\sup_{r>0}\frac{1}{(2r)^n}\int_{Q(x,r)}|f(y)|\dd y.
\]
同一个 $r$ 用于全部坐标。记 $C_n$ 为满足
\begin{equation}\label{eq:weak-definition}
 \lambda\bigl|\{M_nf>\lambda\}\bigr|\le C_n\|f\|_1
 \quad(f\in L^1(\R^n),\ \lambda>0)
\end{equation}
的最小常数。研究固定维数的数值界，需要同时保留共同半径以及一般输入这两个条件：分别估计一维平均的乘积，或只在有限尺度上取极大，都会改变所讨论的问题。

Aldaz 的格点方法给出每个维数的显式下界\cite{Aldaz2000}。令
\[
 u_n=2^{1/n},\quad v_n=(3/2)^{1/n},\quad a_n=(1+u_n)/2,
\]
其 Proposition 1.4 所给的量为
\begin{equation}\label{eq:Aldaz-bound}
 A_n=\frac32a_n^{n-1}
 -\frac12\left[(3a_n-2v_n)^{n-1}
 -\left(\frac{3-u_n}{2}\right)^{n-1}\right].
\end{equation}
特别地 $A_8\approx1.78390524162816$。另一方面，中心 $\ell_\infty$ 球的通用覆盖估计给出 $C_n\le2^n$\cite{AldazKissing2020}，从而 $C_8\le256$。Stein--Str\"omberg 的选择法建立了一般对称凸体的 $O(n\log n)$ 弱型估计\cite{SteinStromberg1983}；在固定维数问题中，还需具体计算覆盖与重叠估计中的常数。Aldaz 的发散定理及后续定量增长结果\cite{Aldaz2011,IakovlevStromberg2013}说明方体常数不能取维数一致的上界，但渐近增长与八维显式数值是两个不同层次的问题。

本文在八维同时构造较大的可实现水平集和较小的普遍覆盖常数。上界的核心是：选择过程留下的未捕获部分位于一个显式方体角帽区域；对坐标壳积分，可从粗重叠估计中扣除始终被捕获的部分。下界的核心是：固定七个间距后，首轴间距对应的完整见证体积连续分段仿射，因此可以按有限端点事件寻找构造，再用准确并集算法认证单个构造。

\begin{theorem}\label{thm:main}
对于上述八维中心轴平行方体极大算子，
\begin{equation}\label{eq:main-bound}
 \boxed{\quad
 \frac{2188811179426669}{10^{15}}<C_8<
 \frac{56627260639}{2000000000}\quad}.
\end{equation}
下界由有限、非负、光滑紧支撑的 $L^1$ 输入实现为一个严格较大的弱型比值；上界对所有 $L^1$ 输入与全部正半径成立。
\end{theorem}

式\eqref{eq:main-bound}的显示形式为
$2.188811179426669<C_8<28.3136303195$。
这里与式\eqref{eq:Aldaz-bound}和通用上界 $256$ 作具体比较，不将这两个比较对象称为固定 Lebesgue 背景下的最优界。计算辅助部分的完整根隔离、并集判据与准确分数作为补充材料给出；主文证明有限判据怎样产生允许的 $L^1$ 输入，并说明它们如何进入定理。

我们还研究一种可用于进一步改进上界的方向势方法。Lin 的正核比较原理\cite{Lin2026}将极大算子的弱型估计归约为具有适当源正性的可积主核。对八维有限方向 Cauchy 过程，Green 势的奇异集已经不只包含方向线：秩不超过三的方向子空间都可能产生不可积的小时间点值。本文证明这一奇异集的准确描述，并以全空间分布方法构造正部势的概率源。由此得到一般输入与全尺度的比较定理，以及控制奇异邻域的面上短时范数估计。这个分析结果与定理\ref{thm:main}的几何上界是两条独立论证。

第二节统一记号和所用标准工具。第三节分别证明几何上界、计算辅助下界与方向势的比较定理。第四节给出强 $L^p$ 推论及十六方向模型的具体计算。附录列出准确证书、验证程序和复现方法。
